\documentclass[12pt]{article}
\usepackage[utf8]{inputenc}
\usepackage[T1]{fontenc}
\usepackage{hyperref}
\title{Mirgan Language Resources for K-12 Curriculum Development}
\author{Compiled by Claude Sonnet 5.5 (claude-sonnet-5-5) for Andr\'as Kornai}
\date{}
\begin{document}
\maketitle

\section{Introduction}
Mirgan (ISO 639-3 zrg; Glottocode mirg1238) is a language of India. Alternate names reported in the sources consulted include Mirgami, Mirkan, Panika and Panka. The Global Recordings Network reports that speakers are bilingual in Odia in Orissa and in Halbi and Hindi in Madhya Pradesh. Joshua Project lists a Mirgan people of about 2,300 in India, mainly in Odisha, with Odia as the main language. The two sources do not fully agree on the language situation. The search found no curriculum, school programme, dictionary or grammar for Mirgan. Very little exists in the sources consulted.

\section{Primary Educational Resources}
No primers, textbooks or school materials for Mirgan were found in the searches.

\section{Linguistic and Reference Materials}
No dictionary, grammar or orthography for Mirgan was found. The Global Recordings Network language page gives basic identification and variant names: \url{https://globalrecordings.net/en/language/zrg}.

\section{Religious and Community Resources}
Joshua Project has a people-group profile for the Mirgan in India: \url{https://joshuaproject.net/people_groups/21770}. Its descriptive text appears generic, so it should be treated with caution.

\section{Digital and Audio Resources}
The Global Recordings Network lists recordings for the Batasuna and Jagdal Pur varieties grouped under Mirgan: \url{https://glorecnet-website.k.globalrecordings.net/en/language/25839} and \url{https://glorecnet-website.k.globalrecordings.net/en/language/25840}. The content of these recordings was not checked.

\section{Curriculum Development Context}
Little documentation exists, and the sources suggest speakers use Odia, Halbi or Hindi in education. A native-language curriculum would first need community consultation and basic documentation work. This dossier does not establish how many people speak the language today.
\end{document}